\section{Conclusion}

Conditional domain leakage in EEG representations can be localized with conditional score-Fisher geometry, and on some representations it can even be partially removed at no task cost --- but neither measuring it nor removing it is, on its own, evidence of a cross-subject generalization benefit. On BCI-IV-2a we find a consistent measurement-to-control gap: in a high-dimensional SPD latent the leakage is redundant and not low-rank removable, and a global penalty removes it only by collapsing the representation; in a compact convolutional latent the same diagnostic deletion removes much of the source-side leakage at no task cost, while a separate collapse-free global-LPC sweep shows that reducing the leakage does not improve target accuracy. Deploying each eraser on the held-out target (source-only fit) confirms this: no eraser --- including the optimal linear eraser LEACE --- improves target accuracy, and the NLL movement is not erasure-specific (same-sign under random removal). This deployment conclusion is not specific to 2a: across four additional real EEG datasets ($129$ further subjects, both backbones), no principled source-fitted eraser --- including optimal linear erasure --- yields a practically meaningful target-accuracy gain in any valid dataset--backbone cell, and on binary datasets naive erasure can even drive the task to chance (\S4.7).

We therefore recommend treating selective conditional invariance as a \textbf{certified intervention with refusal} rather than an always-on regularizer: localize the leakage, certify whether deleting it is both safe for the task and useful, and abstain --- return the identity map --- when it cannot be certified. Framing the decision to \emph{not} intervene as a first-class outcome is, we argue, the honest and practical stance for conditional-invariance methods in EEG.
